\documentclass[11pt,a4paper]{article}

\usepackage[a4paper, margin=2.5cm]{geometry}
\usepackage{amsmath, amssymb, amsthm}
% Engine-agnostic Unicode setup. tectonic (the default builder here)
% runs the XeLaTeX engine, which handles UTF-8 natively via fontspec;
% pdfLaTeX needs the legacy inputenc/fontenc combo. iftex lets the
% same preamble work with both — author names like "Ilhan Karić" then
% render correctly regardless of which engine builds the PDF. Without
% this branch, XeLaTeX silently dropped non-ASCII glyphs (e.g. ć →
% nothing) because the legacy T1 font ec-lmr10 lacks them.
\usepackage{iftex}
\ifPDFTeX
  \usepackage[utf8]{inputenc}
  \usepackage[T1]{fontenc}
  \usepackage{lmodern}
\else
  \usepackage{fontspec}
\fi
\usepackage{hyperref}
\usepackage{enumitem}
\usepackage{xcolor}

\theoremstyle{plain}
\newtheorem{theorem}{Theorem}
\newtheorem{lemma}[theorem]{Lemma}
\newtheorem{proposition}[theorem]{Proposition}
\newtheorem{corollary}[theorem]{Corollary}
\newtheorem{conjecture}[theorem]{Conjecture}

\theoremstyle{definition}
\newtheorem{definition}[theorem]{Definition}
\newtheorem{example}[theorem]{Example}
\newtheorem{notation}[theorem]{Notation}

\theoremstyle{remark}
\newtheorem{remark}[theorem]{Remark}
\newtheorem*{sketch}{Sketch}

\title{Covering systems with moduli of the form $p-1$, $p \geq 5$ prime\\
\large{Erdős Problem \#273}}
\author{Ilhan Karić}
\date{\today}

\begin{document}

\maketitle

\begin{abstract}
\emph{Abstract to be written once the first verified result has
landed; the introduction below states the problem.}
\end{abstract}

\section{Introduction}
\label{sec:intro}

A \emph{covering system} of the integers is a finite collection of
arithmetic progressions
\[
  \{ a_i \pmod{m_i} \}_{i=1}^{k}, \qquad m_i \geq 2,
\]
whose union is $\mathbb{Z}$; equivalently, every integer lies in at
least one residue class $a_i \pmod{m_i}$. Covering systems were
introduced by Erdős in 1950, and a sequence of problems concerning
their structure have driven a substantial body of work in
combinatorial number theory. The present manuscript concerns one of
these problems, recorded as Problem \#273 in the Erdős--Graham survey
\cite{ErdosGraham1980}.

\begin{conjecture}[Erdős \#273]
\label{conj:e273}
Does there exist a covering system $\{ a_i \pmod{m_i} \}_{i=1}^{k}$
of\/ $\mathbb{Z}$ such that each modulus $m_i$ has the form
$m_i = p_i - 1$ for some prime $p_i \geq 5$?
\end{conjecture}

Write $\mathcal{M} := \{ p - 1 : p \text{ prime}, \ p \geq 5 \} =
\{4, 6, 10, 12, 16, 18, 22, 28, 30, 36, 40, 42, 46, 52, 58, 60, \ldots\}$
for the set of allowed moduli. The question is whether $\mathbb{Z}$
admits a finite covering whose moduli are drawn from $\mathcal{M}$.

The restriction $p \geq 5$ is essential. Selfridge observed that
allowing $p = 3$ admits a covering system using divisors of
$360 = \operatorname{lcm}(2,3,4,5,6,8,9,10,12)$, all of which are of
the form $p - 1$ for $p \in \{3, 5, 7\}$. The Selfridge construction
crucially uses the modulus $2 = 3 - 1$. Forbidding $p = 3$
eliminates the modulus~$2$ entirely and changes the problem: any
covering by progressions modulo elements of $\mathcal{M}$ must
account for every parity class without ever using the all-purpose
modulus~$2$, and the second-smallest available modulus is~$4$.

Two pieces of recent machinery bear on the question. Hough
\cite{Hough2015} resolved the longstanding Erdős minimum-modulus
problem by establishing a finite upper bound on the smallest modulus
required in any covering system, ruling out covering systems with all
moduli larger than a fixed constant. Balister, Bollobás, Morris and
Sahasrabudhe \cite{BBMS2022} extended Hough's distortion method to a
general non-existence framework for covering systems whose moduli are
constrained to a multiplicatively structured set. Specializing either
machine to $\mathcal{M}$ is a natural avenue, although neither result
has been adapted to the present restriction in the published
literature.

To the author's knowledge, Conjecture~\ref{conj:e273} is open as of
2026; no partial result---a lower or upper bound on the smallest
required modulus, on the number of progressions, or on the least
common multiple of any putative covering---appears in the published
literature.

The main result of this manuscript is recorded in
Section~\ref{sec:main}, currently in preparation.

\section{Main result}
\label{sec:main}

% [sketch] main result placeholder — to be stated and proved by /solve.
\emph{Statement and proof in preparation.}

\begin{thebibliography}{99}

\bibitem{ErdosGraham1980}
P.~Erdős and R.~L.~Graham,
\emph{Old and new problems and results in combinatorial number theory},
L'Enseignement Mathématique \textbf{28} (1980), Problem~\#273,
p.~24.

\bibitem{Hough2015}
B.~Hough,
\emph{Solution of the minimum modulus problem for covering systems},
Annals of Mathematics \textbf{181} (2015), no.~1, 361--382.
\href{https://arxiv.org/abs/1307.0874}{arXiv:1307.0874}.

\bibitem{BBMS2022}
P.~Balister, B.~Bollobás, R.~Morris, and J.~Sahasrabudhe,
\emph{On the Erdős covering problem: the density of the uncovered set},
Inventiones Mathematicae \textbf{228} (2022), no.~1, 377--414.
\href{https://arxiv.org/abs/1811.03547}{arXiv:1811.03547}.

\end{thebibliography}

\end{document}
